\begin{trapbox}
	\begin{description}[style=nextline, leftmargin=2.8em, labelsep=0.5em, itemsep=0.35em]
		\item[\textbf{\textcolor{CTrap}{易错点一（忽略距离范围）}}] 认为距离$d$可以取任意非负值，未检查$d<R$。
		\item[\textbf{\textcolor{CTrap}{正确理解（距离范围限制）：}}] 球心到平面的距离$d$必须满足$0\le d<R$，否则截面不存在（$d=R$时截面为点，$d>R$时无截面）。
		\item[\textbf{\textcolor{CTrap}{错因分析（忽视定义域）：}}] 只记公式形式，忘记几何条件；$d\ge R$时截面退化或不存在，公式失去意义。
	\end{description}

	\medskip
	\begin{description}[style=nextline, leftmargin=2.8em, labelsep=0.5em, itemsep=0.35em]
		\item[\textbf{\textcolor{CTrap}{易错点二（混淆距离与半径）}}] 在计算中误把$d$当作截面半径$r$，或把$r$当作$d$。
		\item[\textbf{\textcolor{CTrap}{正确理解（直角边有别）：}}] $d$是球心到平面的垂线段长度，$r$是截面圆的半径，二者在直角三角形中都是直角边，但数值不同。
		\item[\textbf{\textcolor{CTrap}{错因分析（空间想象不足）：}}] 未能正确构建直角三角形，导致赋值错误。
	\end{description}

	\medskip
	\begin{description}[style=nextline, leftmargin=2.8em, labelsep=0.5em, itemsep=0.35em]
		\item[\textbf{\textcolor{CTrap}{易错点三（忽视适用范围）}}] 使用逆公式$d=\sqrt{R^{2}-r^{2}}$时，未检查$r\le R$，导致根号内为负。
		\item[\textbf{\textcolor{CTrap}{正确理解（截面半径限制）：}}] 截面半径$r$必须满足$0\le r\le R$（等号仅当$d=0$），逆公式才有效。
		\item[\textbf{\textcolor{CTrap}{错因分析（未注意变量约束）：}}] 忽视几何事实：截面圆半径不超过球半径。
	\end{description}
\end{trapbox}
